\subsection{\texorpdfstring{第~\ref{sec:general-relativity}~节}{广义相对论}}\label{subsec:solutions-general-relativity}

\begin{description}[font=\normalfont,leftmargin=2em,labelwidth=1.5em]
\RGsolution{ex:signature-is-constant-on-pseudo-riemannian-manifold}{signature-is-constant-on-pseudo-riemannian-manifold}
设 $p$ 是 $M^n$ 中的一点. $g_p$ 的符号差就是 $T_pM$ 上二次型 $Q_p(v)=g_p(v,v)$ 的符号差. 在局部坐标 $x:V\subset M\to\mathbb{R}^n$ 中, 这个二次型由矩阵 $A(x)=(g_{ij}(x))$ 表示. 矩阵的各元素都是 $x$ 的连续函数, 所以它的特征值也是 $x$ 的连续函数. 事实上, 这些特征值是 $A(x)$ 的特征多项式的根, 而该多项式的系数是 $A(x)$ 各元素的连续函数. 由于 $A(x)$ 的特征值不可能为零 (因为 $g$ 非退化, 从而 $A(x)$ 可逆), 且 $M$ 连通, 这些特征值的符号不能改变. 结论由此得出.

\RGsolution{ex:trace-reversed-einstein-equation}{trace-reversed-einstein-equation}
\begin{enumerate}[label=(\alph*),ref=\alph*,leftmargin=*]
\item 在任一给定的坐标系中, 有
\[
  \operatorname{Ric}-\frac{S}{2}g=8\pi E
  \quad\Longleftrightarrow\quad R_{\mu\nu}-\frac{S}{2}g_{\mu\nu}=8\pi E_{\mu\nu}
\]
($\mu,\nu=0,1,2,3$). 乘以度量矩阵的逆的分量 $g^{\mu\nu}$ 并对 $\mu,\nu$ 求和, 得
\[
  S-\frac{S}{2}\cdot 4=8\pi\sum_{\mu,\nu=0}^3g^{\mu\nu}E_{\mu\nu}
  \quad\Longleftrightarrow\quad S=-8\pi\sum_{\mu,\nu=0}^3g^{\mu\nu}E_{\mu\nu}.
\]
所以 Einstein 场方程可以写成
\[
  \operatorname{Ric}=8\pi E-4\pi\left(\sum_{\mu,\nu=0}^3g^{\mu\nu}E_{\mu\nu}\right)g=8\pi T.
\]
\item 只需注意, 由于 $\nu$ 是单位类时向量,
\[
  \sum_{\mu,\nu=0}^3 g^{\mu\nu}E_{\mu\nu}=\rho\langle\nu,\nu\rangle=-\rho.
\]
\end{enumerate}

\RGsolution{ex:rotating-frame-in-minkowski-spacetime}{rotating-frame-in-minkowski-spacetime}
\begin{enumerate}[label=(\alph*),ref=\alph*,leftmargin=*]
\item 这是因为 Euclid 度量 $\mathrm{d}x\otimes\mathrm{d}x+\mathrm{d}y\otimes\mathrm{d}y+\mathrm{d}z\otimes\mathrm{d}z$ 在柱坐标 $(r,\theta,z)$ 中的表达式就是 $\mathrm{d}r\otimes\mathrm{d}r+r^2\mathrm{d}\theta\otimes\mathrm{d}\theta+\mathrm{d}z\otimes\mathrm{d}z$.
\item 只需注意
\begin{align*}
  \mathrm{d}\theta\otimes\mathrm{d}\theta
    &=(\mathrm{d}\theta'+\omega\mathrm{d}t)\otimes(\mathrm{d}\theta'+\omega\mathrm{d}t)\\
    &=\mathrm{d}\theta'\otimes\mathrm{d}\theta'+\omega^2\mathrm{d}t\otimes\mathrm{d}t
       +\omega\mathrm{d}t\otimes\mathrm{d}\theta'+\omega\mathrm{d}\theta'\otimes\mathrm{d}t.
\end{align*}
\item 对于 $r<\frac{1}{\omega}$, 在新坐标系 $(t,r,\theta',z)$ 中有
\[
  \left\langle\frac{\partial}{\partial t},\frac{\partial}{\partial t}\right\rangle
  =-1+\omega^2r^2<0,
\]
于是 $(r,\theta',z)$ 不变的曲线都是类时曲线. 沿这些曲线, $r,z$ 不变, 而 $\theta=\theta'+\omega t$, 其中 $\theta'$ 不变. 因此, 这些曲线对应于相对于 $(r,\theta,z)$ 不变的惯性观察者以角速度 $\omega$ 作刚性旋转的观察者.
\item 若在 $U$ 上使用局部坐标 $(t,r,\theta',z)$, 在 $\Sigma$ 上使用 $(r,\theta',z)$, 则投影在局部表示为 $\widehat{\pi}(t,r,\theta',z)=(r,\theta',z)$. 因此, 若
\[
  v=v^r\frac{\partial}{\partial r}+v^{\theta'}\frac{\partial}{\partial\theta'}
    +v^z\frac{\partial}{\partial z},
\]
则
\[
  v^\dagger=v^0\frac{\partial}{\partial t}+v^r\frac{\partial}{\partial r}
    +v^{\theta'}\frac{\partial}{\partial\theta'}+v^z\frac{\partial}{\partial z}
\]
对某个 $v^0\in\mathbb{R}$ 成立. 现在有
\[
  \left\langle\frac{\partial}{\partial t},v^\dagger\right\rangle
  =-(1-\omega^2r^2)v^0+\omega r^2v^{\theta'},
\]
所以必有
\[
  v^0=\frac{\omega r^2}{1-\omega^2r^2}v^{\theta'}.
\]
因此
\begin{align*}
  \langle v^\dagger,w^\dagger\rangle
     &=-(1-\omega^2r^2)v^0w^0+\omega r^2\bigl(v^0w^{\theta'}+v^{\theta'}w^0\bigr)\\
     &\quad+v^rw^r+r^2v^{\theta'}w^{\theta'}+v^zw^z\\
     &=v^rw^r+\frac{r^2}{1-\omega^2r^2}v^{\theta'}w^{\theta'}+v^zw^z.
\end{align*}
注意, 这与分别投影到 $v,w$ 的向量 $v^\dagger,w^\dagger$ 的选取无关.
\item 由 (d) 可知, 向量 $w$ 与 $\frac{\partial}{\partial t}$ 正交的条件为
\[
  w^0=\frac{\omega r^2}{1-\omega^2r^2}w^{\theta'}
  \quad\Longleftrightarrow\quad \mathrm{d}t(w)=\frac{\omega r^2}{1-\omega^2r^2}\mathrm{d}\theta'(w)
  \quad\Longleftrightarrow\quad\alpha(w)=0.
\]
因此, 曲线 $c$ 由同时事件组成, 当且仅当 $\alpha(\dot c)=0$. 若 $\gamma:[0,1]\to\Sigma$ 为闭曲线, 且 $c:[0,1]\to U$ 是由同时事件组成并满足 $\pi\circ c=\gamma$ 的曲线, 则
\begin{align*}
  t(c(1))-t(c(0))
    &=\int_0^1\mathrm{d}t(\dot c)
      =\int_0^1\frac{\omega r^2}{1-\omega^2r^2}\mathrm{d}\theta'(\dot c)\\
    &=\int_0^1\frac{\omega r^2}{1-\omega^2r^2}\mathrm{d}\theta'(\dot\gamma).
\end{align*}
由于
\[
  \mathrm{d}\left(\frac{\omega r^2}{1-\omega^2r^2}\mathrm{d}\theta'\right)
  =\frac{2\omega r}{(1-\omega^2r^2)^2}\mathrm{d}r\wedge\mathrm{d}\theta'\ne 0,
\]
一般而言 $t(c(1))\ne t(c(0))$. 所以, 若旋转观察者沿一条闭路径同步各个时钟, 他们会发现最初的时钟与其自身并不同步.
\end{enumerate}
\end{description}
